\documentclass[11pt]{article}
\usepackage[margin=2.2cm]{geometry}
\usepackage{amsmath,amssymb}
\usepackage{longtable}
\newcommand{\Tr}{\operatorname{Tr}}
\newcommand{\STr}{\operatorname{STr}}
\newcommand{\dm}{\dot{-}}
\newcommand{\dpl}{\dot{+}}
\newcommand{\pb}{\bar{\partial}}
\newcommand{\cI}{\mathcal{I}}
\newcommand{\cO}{\mathcal{O}}
\newcommand{\bQ}{\mathbf{Q}}
\setlength{\parindent}{0pt}
\setlength{\parskip}{4pt}
\begin{document}

\section*{Reading notes: arXiv:2306.01039}

\textbf{Paper.} K.~Budzik, D.~Gaiotto, J.~Kulp, B.~R.~Williams, J.~Wu, M.~Yu, \emph{Semi-Chiral Operators in 4d $\mathcal N=1$ Gauge Theories}, arXiv:2306.01039 [hep-th] (id verified on the abs page and on the PDF title page); v1 1 Jun 2023, v2 3 Sep 2024 (= published version), JHEP 05 (2024) 245, DOI 10.1007/JHEP05(2024)245. 55+20 pages.

\textbf{Source used.} PDF of v2 (arxiv.org/pdf/2306.01039v2, read page by page). Equation numbers below are those of the v2 PDF (section-numbered, e.g.\ (3.40)). The arXiv HTML rendering of v2 numbers equations consecutively and splits multi-line displays; key correspondences: (3.24)$=$HTML(55), (3.27)$=$(58), (3.28)$=$(59), (3.30)$=$(61), (3.33)$=$(64), (3.36)$=$(67), (3.39)$=$(70), (3.40)$=$(71), (3.41)--(3.45)$=$(72)--(76), (4.1)$=$(77), (5.20)$=$(128), (5.21)$=$(129). Page numbers ``p.\,$n$'' are the printed page numbers (PDF page $=n+2$).

\textbf{Name used in these notes.} ``$Q_1$'' always means the paper's one-loop term of the loop-corrected twisted differential, (3.28), (D.13). It is \emph{not} the spinor parameter $Q_1\in\Sigma$ of the manuscript (Prop.~prop:c7-susy). Prose of the paper is paraphrased; displayed formulas are transcribed.

\section{Conventions of the paper}

\textbf{Signature, spinors.} Euclidean $\mathcal N=1$ algebra; $\mathrm{Spin}(4)\cong SU(2)\times SU(2)$ (p.\,5). Footnote~7: the paper \emph{switches} the standard (Wess--Bagger [34]) assignment of dotted/undotted indices.
\begin{align}
&\{Q_{\dot\alpha},\bar Q_\beta\}=P_{\dot\alpha\beta}, \tag{2.1}\\
&\{Q_{\dot\alpha},Q_{\dot\beta}\}=\{\bar Q_\alpha,\bar Q_\beta\}=0 . \tag{2.2}
\end{align}
Rotations: $M_{\dot\alpha\dot\beta}$ and $\bar M_{\alpha\beta}$ (p.\,5). No normalisation of $P_{\dot\alpha\beta}$ in terms of $\sigma$-matrices is given.

\textbf{Twisting supercharge and complex structure} (p.\,6). $Q\equiv Q_{\dm}$; $z^\alpha\equiv x^{\dpl\alpha}$ holomorphic, $\bar z^\alpha\equiv x^{\dm\alpha}$ antiholomorphic; footnote~9: $(z_\alpha)^*=\epsilon_{\alpha\beta}\bar z^\beta$.
\begin{equation}
\{Q,\bar Q_\alpha\}=\partial_{\dm\alpha}=\partial_{\bar z^\alpha} . \tag{2.3}
\end{equation}
Surviving rotations: the $SU(2)$ generated by $\bar M_{\alpha\beta}$ (holomorphic symplectic rotations, preserving $\mathrm d^2z=\mathrm dz^1\wedge\mathrm dz^2$). With $U(1)_R$: $M_R\equiv M_{\dpl\dm}-R$ extends it to $U(2)$. Footnote~10: $[M_{\dpl\dm},Q_{\dm}]=-Q_{\dm}$, $[R,Q_{\dm}]=-Q_{\dm}$. Cohomological degree $C\equiv \mathrm{gh}-M_{\dpl\dm}$.

\textbf{Charges} (Table~1, p.\,50):
\[
\begin{array}{c|ccccc}
 & R & M_{\dpl\dm} & M_R & C_R & C\\\hline
Q & -1 & -1 & 0 & 1 & 1\\
b & 0 & 2 & 2 & 0 & -2\\
c & -1 & -1 & 0 & 1 & 1\\
\partial_\alpha & 0 & 1 & 1 & 0 & -1
\end{array}
\qquad M_R=M_{\dpl\dm}-R,\quad C_R=\mathrm{gh}-R,\quad C=C_R-M_R=\mathrm{gh}-M_{\dpl\dm}.
\]

\textbf{Superconformal relations} (for SCFTs, kinematics only):
\begin{align}
&\{Q_{\dm},S_{\dpl}\}=\Delta+\tfrac32R-M_{\dpl\dm}, \tag{2.4}\\
&\{Q_{\dot\alpha},S_{\dot\beta}\}=\big(\Delta+\tfrac32R\big)\epsilon_{\dot\alpha\dot\beta}-M_{\dot\alpha\dot\beta}. \tag{2.5}
\end{align}
Short $\mathfrak{su}(1|1)$ multiplets: $\Delta=M_{\dpl\dm}-\frac32R$ (p.\,7). Statement (p.\,7): in SCFTs $Q$-cohomology classes are annihilated by $M_{\dpl\dpl}$ and have non-positive cohomological degree; degree $0$ = chiral ring.

\textbf{Chiral superspace and descent} (p.\,8--9). $\bar Q_\alpha=-\partial_{\bar\theta^\alpha}$ on superfields; $\bar\theta^\alpha\leftrightarrow\mathrm d\bar z^\alpha$.
\begin{align}
&\pb\equiv\bar\theta^\alpha\partial_{\bar z^\alpha}, \tag{2.6}\\
&\cO[\bar\theta]=e^{\bar\theta^\alpha\bar Q_\alpha}\cO^{(0)}=\cO^{(0)}+\cO^{(1)}+\cO^{(2)}, \tag{2.7}\\
&\{Q+\pb,\bar Q_\alpha\}=0, \tag{2.8}\\
&(Q+\pb)\cO[\bar\theta]=(Q\cO)[\bar\theta], \tag{2.9}\\
&Q\cO^{(k)}+\pb\cO^{(k-1)}=(Q\cO)^{(k)}, \tag{2.10}\\
&D_{\dm}\cO=(Q+\pb)\cO=0\quad\text{(semi-chiral)}, \tag{2.11}\\
&Q\cO^{(1)}+\pb\cO^{(0)}=0,\qquad Q\cO^{(2)}+\pb\cO^{(1)}=0, \tag{2.12}\\
&\cO=(Q+\pb)\,\mathcal U\quad\text{if }\cO^{(0)}=Q\,\mathcal U^{(0)}. \tag{2.13}
\end{align}
$Q+\pb$ coincides with the superspace derivative $D_{\dm}$ of the physical theory (p.\,8).

\textbf{$\lambda$-bracket} ($n$ complex dimensions):
\begin{equation}
\{\cO_1\,{}_\lambda\,\cO_2\}=\oint_{S^{2n-1}}e^{\lambda\cdot z}\,\mathrm d^nz\,\cO_1(z)\cO_2(0). \tag{2.18}
\end{equation}

\textbf{Lie algebra} (p.\,19--20). Killing form $\delta^{AB}$, adjoint indices raised/lowered freely; trace in the defining representation;
\begin{align}
&\Tr t^At^B=-\tfrac12\delta^{AB},\qquad \Tr t^A[t^B,t^C]=-\tfrac12f^{ABC},\qquad [t^A,t^B]=f^{ABC}t^C, \tag{3.20}\\
&\textstyle\sum_{A,B}f^{ABC}f^{ABD}=N\delta^{CD}, \tag{3.21}\\
&SU(2):\ t^A=-\tfrac i2\sigma^A,\ f^{ABC}=\epsilon^{ABC};\qquad Qb^A=\epsilon^{ABC}c^Bb^C,\quad Qc^A=\tfrac12\epsilon^{ABC}c^Bc^C, \tag{3.22}\\
&b^2=b_Ab^A=-2\Tr b^2. \tag{3.23}
\end{align}

\textbf{Tree bracket} (App.~B, p.\,57).
\begin{align}
&\{b^A,c^B\}=\delta^{AB}, \tag{B.1}\\
&\{\cO_1,\cO_2\}=\delta^{AB}\frac{\delta\cO_1}{\delta c^A}\frac{\delta\cO_2}{\delta b^B}+(-1)^{|\cO_1|}\delta^{AB}\frac{\delta\cO_1}{\delta b^A}\frac{\delta\cO_2}{\delta c^B}, \tag{B.2}\\
&\frac{\delta}{\delta c^A}\big(c^Bc^C\big)=\delta^{AB}c^C-\delta^{AC}c^B. \tag{B.3}
\end{align}
Footnote~54 (App.~B.2): $\psi^\alpha\chi_\alpha=\psi_2\chi_1-\psi_1\chi_2=-\psi_\alpha\chi^\alpha$.

\textbf{Coupling and $\hbar$.} The BF action (3.2) contains no coupling; $\tau$ appears only in (3.4), a total derivative in BF theory (instantons only). $\hbar$ appears in (3.27) and in (D.13): $Q=Q_0+\hbar Q_1+\hbar^2Q_2+\cdots$; (3.28) writes $\bQ=Q_0+Q_1+\cdots$ with $Q_n$ the $n$-loop term. No relation between $\hbar$ and the physical $g^2$ (or $g^2N/16\pi^2$) is given.

\section{Twisted theory and its relation to the untwisted fields}

\textbf{Pure $\mathcal N=1$ gauge theory = holomorphic BF theory} (p.\,13--15). The identification (in BV, refs.\ [22,41,50]) changes the field content but not the $Q$-cohomology. $Y=\mathbb C^2$.
\begin{align}
&b=b^{(0)}+b^{(1)}+b^{(2)}\in\Omega^{0,\bullet}(\mathbb C^2,\mathfrak g^\vee),\qquad c=c^{(0)}+c^{(1)}+c^{(2)}\in\Omega^{0,\bullet}(\mathbb C^2,\mathfrak g)[1], \tag{3.1}\\
&\int\mathrm d^2z\,\Tr\,b\Big(\pb c-\tfrac12[c,c]\Big), \tag{3.2}\\
&\int\mathrm d^2z\,\Tr\,b^{(0)}\Big(\pb c^{(1)}-\tfrac12[c^{(1)},c^{(1)}]\Big), \tag{3.3}\\
&\int\mathrm d^2z\,\tau\,\Tr\,\partial_\alpha c\,\partial^\alpha c . \tag{3.4}
\end{align}
Degrees: $c^{(0)}$ ghost, degree $1$; $c^{(1)}$ degree $0$ = antiholomorphic part of the connection; $b^{(0)}$ degree $-2$; the holomorphic part of the connection is absent; the 2-form part of $\pb c-\frac12[c,c]$ is the antiholomorphic field strength. Footnote~21: $b^{(i)}$ is naturally a $(2,i)$ form; $\mathrm d^2z$ fixed, so $(0,i)$. Footnote~22: $c$ adjoint, $b$ coadjoint, identified by an invariant pairing.

Dictionary to the physical fields (p.\,14, stated up to field redefinitions, at the level of $Q$-cohomology of the free theory):
\begin{itemize}
\item $(2,0)$ part of $F_{\dpl\dpl}$ $\leftrightarrow$ $\mathrm d^2z\,b^{(0)}$;
\item $\bar\lambda_\alpha$ ($\alpha=1,2$), with $\epsilon_{\alpha\beta}\partial^\alpha\bar\lambda^\beta=\partial_\alpha\bar\lambda^\alpha=0$ $\leftrightarrow$ $\partial_{z^\alpha}c^{(0)}$; superfield: $\overline{\mathcal W}_\alpha$ (anti-chiral) $\leftrightarrow$ $\partial_{z^\alpha}c$;
\item kinetic F-term $\int\mathrm d^2\bar\theta\,\tau\,\overline{\mathcal W}_\alpha\overline{\mathcal W}^\alpha$ $\leftrightarrow$ (3.4).
\end{itemize}
The basic action corresponds to a self-dual limit of the twisted theory (p.\,14). Footnote~20: on a curved complex surface $Y$ the one-loop quantisation has an anomaly $\propto c_1(T_Y)\times$ adjoint quadratic Casimir.

\textbf{Chiral multiplet} (§2.5):
\begin{align}
&\gamma\equiv\Phi, \tag{2.23}\\
&\beta\equiv D_{\dpl}\bar\Phi, \tag{2.24}\\
&S_\alpha=\beta\partial_\alpha\gamma-\Delta\,\partial_\alpha(\beta\gamma), \tag{2.25}\\
&\partial_\alpha S^\alpha=\partial_\alpha\beta\,\partial^\alpha\gamma, \tag{2.26}\\
&(Q+\pb)\beta=\partial_\gamma\mathcal W(\gamma), \tag{2.27}\\
&(Q+\pb)\,\beta\partial_\alpha\gamma=\partial_\alpha\big[\mathcal W-\Delta\,\gamma\partial_\gamma\mathcal W\big]. \tag{2.28}
\end{align}
$J=\beta\gamma$ from $|\Phi|^2$; $\Delta$ in (2.25) $\propto$ R-charge of $\Phi$.

\textbf{Matter, tree level} (§4.1). Chirals in $R$: bosonic $\gamma$ in $R$, fermionic $\beta$ in $\bar R\simeq R^*$.
\begin{align}
&\int\mathrm d^2z\,\beta\big(\pb\gamma-[c,\gamma]\big)+W(\gamma), \tag{4.1}\\
&Q_0b=-\pb b+[c,b]+\gamma\beta,\qquad Q_0\gamma=-\pb\gamma+[c,\gamma],\qquad Q_0\beta=-\pb\beta+[c,\beta]+\partial_\gamma W, \tag{4.2}\\
&S_\alpha=\Tr b\partial_\alpha c+c_1\tilde\Gamma^A\partial\,\Gamma_A+c_2\,\partial\tilde\Gamma^A\Gamma_A\quad(\text{SQCD, constants }c_1,c_2). \tag{4.3}
\end{align}
With matter: $\Tr b^2$ is not tree-level $Q$-closed; in the SQCD discussion the stress tensor and its derivatives are stated not to be exact, at least perturbatively (p.\,28). Chiral operators $\leftrightarrow$ $G$-invariant polynomials in $\gamma$ modulo multiples of $\partial_\gamma W$.

\textbf{$\mathcal N=4$ SYM} (p.\,28, 31). Twist = holomorphic $bc$ system in $\mathfrak{gl}(N)$ + three adjoint $\beta\gamma$ systems $(\beta_i,\gamma^i)$, $i=1,2,3$; superpotential $\Tr[\gamma^1,\gamma^2]\gamma^3$ (p.\,28), written $W=\Tr\gamma^1[\gamma^2,\gamma^3]$ on p.\,31 (equal by cyclicity). Following [19]:
\begin{align}
&C(\varepsilon_1,\varepsilon_2,\varepsilon_3)\equiv c+\varepsilon_i\gamma^i-\tfrac12\epsilon^{ijk}\varepsilon_i\varepsilon_j\beta_k-\varepsilon_1\varepsilon_2\varepsilon_3\,b, \tag{4.4}\\
&Q_0C(\varepsilon_i)=\tfrac12[C(\varepsilon_i),C(\varepsilon_i)], \tag{4.5}\\
&C^\bullet\big(\mathfrak{gl}_N\otimes\mathbb C[[z_1,z_2,\varepsilon_i]]\,|\,\mathfrak{gl}_N\big), \tag{4.6}
\end{align}
\begin{equation}
Q_0c=\tfrac12[c,c],\quad Q_0b=[c,b]+\sum_{i=1}^3[\beta_i,\gamma^i],\quad Q_0\gamma^i=[c,\gamma^i],\quad Q_0\beta_i=[c,\beta_i]+\tfrac12\sum_{j,k=1}^3\epsilon_{ijk}[\gamma^j,\gamma^k]. \tag{4.14}
\end{equation}
Enhanced tree-level symmetry: holomorphic vector fields on $\mathbb C^{2|3}$ (coordinates $z_\alpha,\varepsilon_i$) preserving $\mathrm d^2z\,\mathrm d^3\varepsilon$ (p.\,29). \emph{No one-loop formula for $\mathcal N=4$ is given in this paper.}

\section{Local operators: free and tree level}

\begin{align}
&\mathcal V\equiv\mathbb C\Big[\partial^n_{z^1}\partial^m_{z^2}c^{(0)},\ \partial^k_{z^1}\partial^l_{z^2}b^{(0)},\ n+m>0,\ k+l\geq0\Big]^G, \tag{3.5}\\
&Qc=-\pb c+\tfrac12[c,c],\qquad Qb=-\pb b+[c,b], \tag{3.6}\\
&(Q+\pb)\partial_\alpha c=[c,\partial_\alpha c],\quad (Q+\pb)\partial_\alpha\partial_\beta c=[c,\partial_\alpha\partial_\beta c]+[\partial_\alpha c,\partial_\beta c],\quad (Q+\pb)\partial_\alpha b=[c,\partial_\alpha b]+[\partial_\alpha c,b], \tag{3.7}\\
&Q=Q_{\rm free}+Q_{\rm int},\qquad Q_{\rm free}=-\pb, \tag{3.8}\\
&Q_{\rm int}c=\tfrac12[c,c],\qquad Q_{\rm int}b=[c,b]. \tag{3.9}
\end{align}
$c^{(0)}$ itself is excluded from local operators (only derivatives; constant gauge transformations imposed as strict invariance, p.\,15). Free theory: kinetic term pairs $b^{(0)}$ with $c^{(1)}$ and $b^{(1)}$ with $c^{(0)}$, so $b^{(0)},c^{(0)}$ have regular OPE and their composites have no renormalisation ambiguity in the free theory; $\Delta_{\rm BV}$ must be regularised (footnote~24); quantum corrections do not change $\mathcal V$ in the free theory (p.\,15--16). Tree level: $Q_{\rm int}$ reproduces the categorification literature [18,19,51], with $Q_{\dm}$ acting on $D_{z^\alpha}$ as $[\bar\lambda_\alpha,\cdot\,]$ (p.\,17).
\begin{align}
&C^\bullet\big(\mathfrak g[[z_1,z_2]]\,|\,\mathfrak g;\ \mathrm{Sym}(\mathfrak g[[z_1,z_2]]^*)\big), \tag{3.10}\\
&A_n=\Tr b^n,\ n\geq1, \tag{3.11}\\
&C(\varepsilon)\equiv c+\varepsilon b,\qquad Q_{\rm int}C(\varepsilon)=\tfrac12[C(\varepsilon),C(\varepsilon)], \tag{3.12}\\
&\xi=\partial_\varepsilon:\ c\to b,\qquad \tilde\xi_\alpha=\varepsilon\partial_\alpha:\ b\to\partial_\alpha c, \tag{3.14}\\
&B_{a-1,\alpha}\equiv\tilde\xi_\alpha A_a,\qquad C_{a-2}\equiv\tilde\xi_1\tilde\xi_2A_a, \tag{3.15}\\
&\xi B_{a-1,\alpha}=\partial_\alpha A_a,\qquad \xi C_{a-2}=2\partial_\alpha B^\alpha_{a-1}, \tag{3.16}\\
&\partial_\alpha B^\alpha_{a-1}=\tilde\xi^\alpha\partial_\alpha A_a, \tag{3.17}\\
&A_n=\Tr b^n\ (n\geq2),\qquad B_{n,\alpha}=\STr b^n\partial_\alpha c\ (n\geq1),\qquad C_n=\STr b^n\partial_\alpha c\,\partial^\alpha c\ (n\geq0), \tag{3.18}\\
&S_\alpha=\Tr b\,\partial_\alpha c. \tag{3.19}
\end{align}
$\xi$ generated by $A_2=\Tr b^2$; $\tilde\xi_\alpha$ by $\Tr c\,\partial_\alpha c$ (primitive of $C_0$) (p.\,19). Large $N$ (§5): $\Tr[\mathrm dC]^n$ with $\mathrm dC=\mathrm d\varepsilon\,b+\mathrm dz_\alpha\partial_\alpha c+\varepsilon\,\mathrm dz_\alpha\partial_\alpha b$ (5.8);
\begin{align}
\Tr[\mathrm dC]^n&=\mathrm d\varepsilon^nA_n+\mathrm dz_\alpha\mathrm d\varepsilon^{n-1}B_{n-1,\alpha}+\mathrm dz^1\mathrm dz^2\mathrm d\varepsilon^{n-2}C_{n-2}\nonumber\\
&\quad+\varepsilon\,\mathrm dz_\alpha\mathrm d\varepsilon^{n-1}n^{-1}\partial_\alpha A_n+\varepsilon\,\mathrm dz^1\mathrm dz^2\mathrm d\varepsilon^{n-2}(n-1)^{-1}\partial_\alpha B_{n-1,\alpha}, \tag{5.10}\\
&A_n=\Tr b^n\ (n\geq1),\quad B_{n,\alpha}=\Tr b^n\partial_\alpha c\ (n\geq1),\quad C_n=\Tr b^n\partial_\alpha c\,\partial^\alpha c\ (n\geq0). \tag{5.19}
\end{align}
Large-$N$ single-trace tree cohomology: (5.19) with derivatives, plus derivatives of $\Tr c$ (p.\,41). B.2 normalisation (different): $A_n=\frac{(2i)^n}{n}\Tr b^n$, $A_n[b+\eta^\alpha\partial_\alpha c]=A_n+\eta^\alpha B_{n-1,\alpha}+\eta^1\eta^2C_{n-2}$ (B.12); tree $\lambda$-brackets (B.13) hold up to $Q_0$-exact terms. App.~B.1: stress tensor of the $bc$ system written $2\Tr b\,\partial_\alpha c$ (factor $2$ relative to (3.19)), with
\begin{align}
&g_{m,n}=z_1^{m+1}z_2^{n+1}\{\partial_\alpha S^\alpha,-\},\qquad \partial_\alpha S^\alpha=2\Tr\partial_\alpha b\,\partial^\alpha c=-\partial_\alpha b_A\partial^\alpha c^A, \tag{B.8}\\
&g_{m,n}\cO=-X_{f_{m,n}}\cO,\qquad X_f=\frac{\partial f}{\partial z_1}\frac{\partial}{\partial z_2}-\frac{\partial f}{\partial z_2}\frac{\partial}{\partial z_1}, \tag{B.10, B.4}\\
&g_{-1,0}\cO=\partial_1\cO,\quad g_{0,-1}\cO=-\partial_2\cO,\quad -\tfrac12(g_{1,-1}+g_{-1,1})\cO=(x_1\partial_2-x_2\partial_1)\cO. \tag{B.11}
\end{align}

\section{Definition of the loop-corrected differential (what is defined)}

\begin{align}
&\pb+\hbar\Delta_{\rm BV}+\{\cI,\cdot\}_{\rm BV}, \tag{3.27}\\
&\bQ=Q_0+Q_1+\cdots. \tag{3.28}
\end{align}
(3.27): full quantum differential on operators; proper definition needs renormalisation of UV and IR divergences (p.\,21). Conjecture (p.\,21): the cohomology of (3.27) is computed by a differential $\bQ$ on $\mathcal V$, a quantum-corrected $Q_{\rm int}$. Operational definition (p.\,21--22): regularise $\cO\,e^{\int\cI}$, act with $Q_{\rm free}$, rewrite the result as the regularisation of $(\bQ\cO)\,e^{\int\cI}$. $Q_n$: $n$-loop diagrams of the holomorphic theory. Footnote~29: cohomology of the formal series via $Q_0^2=0$ and homotopy transfer (App.~D).
\begin{align}
&\int_{\mathbb R^4}\cI^{(2)}(z,\bar z)\,\mathrm d^2z,\qquad \cI=-\tfrac12\Tr b[c,c], \tag{3.29}\\
&Q_0\cO=Q_{\rm free}\Big[\int_{|z|>\epsilon}\cI(z,\bar z)\mathrm d^2z\ \cO(0,0)\Big]=\int_{|z|=\epsilon}\cI(z,\bar z)\mathrm d^2z\,\cO(0,0), \tag{3.30}\\
&Q_0\cO=\{\cI,\cO\}\quad(\lambda=0\text{ specialisation of (2.18)}), \tag{3.31}\\
&\int_{|z|=\epsilon}\mathrm d^2z\,F(z,\bar z)P(z,\bar z)\to F(0,0)+O(\epsilon), \tag{3.32}\\
&Q_n\cO=\frac1{(n+1)!}\{\cI,\cI,\cdots,\cI,\cO\}, \tag{3.33}\\
&\tfrac12\{\cI,\cI\}+\tfrac16\{\cI,\cI,\cI\}+\cdots=0, \tag{3.34}\\
&\int_{\mathbb R^{4n}}\prod_{i=1}^n\cI(z_i,\bar z_i)\mathrm d^2z_i\,\cO(0,0), \tag{3.35}\\
&P(x,\bar x,\mathrm d\bar x)\equiv\frac1{4\pi^2}\frac{\bar x^2\mathrm d\bar x^1-\bar x^1\mathrm d\bar x^2}{|x|^4},\qquad |x|^2=x^1\bar x^1+x^2\bar x^2, \tag{3.36}\\
&\pb P(x,\bar x,\mathrm d\bar x)=\mathrm d\bar x^1\mathrm d\bar x^2\,\delta^4(x),\qquad \pb=\mathrm d\bar x^1\tfrac{\partial}{\partial\bar x^1}+\mathrm d\bar x^2\tfrac{\partial}{\partial\bar x^2}. \tag{3.37}
\end{align}
(3.30): regulator for $Q_0$ = interaction switched off in a ball of radius $\epsilon$ around $\cO$ (BRST anomaly at the boundary); as an element of $Q_{\rm free}$-cohomology it is $\epsilon$-independent. (3.38): each propagator $P(x-y)$ contracts $c(x)$ with $b(y)$ and removes one $b$ and one $c$. Each $\cI$ adds two $c$ and one $b$. $Q_n$: integral over $2n+2$ antiholomorphic variables, hence $2n+1$ propagators; net number of $c$ increases by $1$, of $b$ decreases by $n$; each type of derivative increases by $n$ (p.\,24). Higher brackets: quadratic axioms + MC condition (3.34) $\Rightarrow$ $\bQ$ nilpotent order by order (argued; details in [26]).

Procedure (p.\,23--24), paraphrased: (1) integral (3.35); (2) maximal number of Wick contractions with \emph{regularised} propagators; (3) act with $Q_{\rm free}$, integrate $\pb$ by parts onto the propagators; (4) set descendants to zero, project to the origin ($\bar z_i=0$, Taylor in $z_i$) to get an element of $\mathcal V$; (5) the resulting Feynman integrals are claimed finite and regulator independent. The regulator of the propagators is not specified in this paper.

Tree level (p.\,24--25): single internal line; Leibniz rule; commutes with $\partial_\alpha$;
\begin{equation}
Q_0b(x)=\{\cI,b\}(x)=\int_{|y-x|>\epsilon}\mathrm d^2y\,\pb_yP(y-x)[c(y),b(y)]=[c(x),b(x)]-\pb_x\int_{|y-x|>\epsilon}\mathrm d^2y\,P(y-x)[c(y),b(y)], \tag{3.39}
\end{equation}
i.e.\ $[c,b]$ up to a $Q_{\rm free}$-exact term; $Q_0c=\{\cI,c\}=\frac12[c,c]$; $Q_0=Q_{\rm int}$.

Nilpotency and filtration (p.\,26--27):
\begin{align}
&\bQ^2\cO=0\quad\forall\cO\in\mathcal V, \tag{3.41}\\
&\sum_{n_1+n_2=n}Q_{n_1}Q_{n_2}\cO=0\quad\forall\cO\in\mathcal V,\ \forall n\geq0, \tag{3.42}\\
&Q_nV(N_b,N_c,J_1,J_2)=0\quad\text{if }N_b+N_c<n+1, \tag{3.43}\\
&Q_n:V(n+1,0,J_1,J_2)\to V(1,1,J_1+n,J_2+n), \tag{3.44}\\
&Q_n:V(n,1,J_1,J_2)\to V(0,2,J_1+n,J_2+n). \tag{3.45}
\end{align}
$V(N_b,N_c,J_1,J_2)$: monomials with $N_b$ letters $b$, $N_c$ letters $c$ (with derivatives), $J_i$ derivatives in $z_i$. $Q_n$ linear; derivation on monomials longer than $n+1$. First non-trivial equation of (3.42): $\{Q_0,Q_1\}=0$. $Q_1^2\neq0$ (observed), hence $Q_2\neq0$; [32] computes the only two-loop diagram, non-vanishing. Auxiliary $U(1)_b$ ($b$ charge $1$, $c$ charge $0$): $Q_0$ preserves it, $Q_n$ lowers it; the filtration makes $\bQ$ well defined on monomials and the spectral sequence well posed.

Homotopy transfer (App.~D.2):
\begin{align}
&Q=Q_0+\hbar Q_1+\hbar^2Q_2+\cdots, \tag{D.13}\\
&\Delta\equiv\iota_0\pi_0+\{h_0,Q\}=1_X+\hbar\{h_0,Q_1\}+\hbar^2\{h_0,Q_2\}+\cdots, \tag{D.14}\\
&\tilde Q=\pi_0(\hbar Q_1+\cdots)\Delta^{-1}i_0=\hbar Q^{(1)}_0+\hbar^2Q^{(1)}_1+\cdots. \tag{D.15}
\end{align}
Iterating gives a limiting $U_0^{(\infty)}$, taken as the definition of perturbative cohomology (requires a filtration).

\section{One-loop formulas (what is computed)}

\textbf{Structure} (§3.2.3, App.~C).
\begin{equation}
Q_1\cO\equiv\tfrac12\{\cI,\cI,\cO\}\quad(\text{all }\lambda=0). \tag{C.1}
\end{equation}
Only triangle diagrams: two fields of $\cO$ and two fields of each of the two vertices $\cI(x),\cI(y)$ contracted. On a single field the triple bracket vanishes identically. On longer monomials, sum over pairs:
\begin{equation}
Q_1\big(f_1\dots f_{i-1}f_if_{i+1}\dots f_{j-1}f_jf_{j+1}\dots\big)=\sum_{i<j}(-1)^{\cdots}Q_1(f_if_j)\big(f_1\dots f_{i-1}f_{i+1}\dots f_{j-1}f_{j+1}\dots\big), \tag{C.2}
\end{equation}
signs from moving $f_i,f_j$ to the front. Properties stated (p.\,25): respects gauge symmetry and the perturbative $U(2)$; turns one $b$ into a $c$ and adds two holomorphic derivatives with contracted spacetime indices.

\textbf{Pairs, from the diagram computed in [32]} (App.~C; open adjoint indices; the $\delta_{AB}$ structure is that of $SU(2)$, see §8):
\begin{align}
&Q_1(b_Ab_B)=\partial_\alpha b_A(x)\partial^\alpha c^B(x)+\partial_\alpha b_B(x)\partial^\alpha c^A(x)-2\delta_{AB}\,\partial_\alpha b_C(x)\partial^\alpha c^C(x), \tag{C.3}\\
&Q_1(b_Ac^B)=(-1)\Big(\partial_\alpha c^A(x)\partial^\alpha c^B(x)-\delta^B_A\,\partial_\alpha c_D(x)\partial^\alpha c^D(x)\Big). \tag{C.4}
\end{align}

\textbf{$\Tr b^2$} (p.\,25):
\begin{equation}
Q_1\Tr b^2\propto\partial_\alpha\big(\Tr b\,\partial^\alpha c\big)=\partial_\alpha S^\alpha . \tag{3.40}
\end{equation}
Claims: $\partial_\alpha S^\alpha$, non-trivial at tree level, is $Q$-exact at one loop; by (3.43) all higher $Q_n$ act trivially on $\Tr b^2$, so (3.40) receives no further perturbative correction. No coefficient is printed in (3.40).

\textbf{Planar, single-trace} (§5.2). Normalisation: $N^{-1}$ per propagator, $N$ per interaction vertex, single traces with an overall power of $N$ (not specified); the planar differential maps single-trace to single-trace with no $N$; $1\to2$ traces suppressed by $N^2$, $2\to1$ of order $1$.
\begin{align}
&Q_1A_n\sim n\Tr b^{n-2}\partial_\alpha b\,\partial^\alpha c, \tag{5.20}\\
&Q_1A_n=\frac{2n}{n-1}\,\partial_\alpha B^\alpha_{n-1}. \tag{5.21}
\end{align}
(5.21) after reordering letters with tree-exact terms such as $Q_0\Tr b^{k_1}\partial_\alpha b\,b^{k_2}\partial^\alpha b$, i.e.\ on $Q_0$-cohomology. Further statements (p.\,42--43): $Q_1B_{n,\alpha}$ must be a multiple of $\partial_\alpha C_n$ up to tree-exact terms unless an accidental cancellation occurs (coefficient not given); $Q_1C_n=0$ (no single-trace tree classes with three $c$); one-loop cohomology: $U(N)$ keeps $C_n$ (no derivatives) and derivatives of $\Tr c$; $SU(N)$ keeps only $C_n$; derivatives of $C_n$ trivial; footnote~42: the $2\to1$ part does not change this; quantum numbers forbid further perturbative corrections; even ghost number forbids non-perturbative ones. Index check: $x\to pq$ once loop corrections are included (p.\,49, (7.7)).

\textbf{Bootstrap of $Q_1$ for $SU(2)$} (App.~C). Ansatz
\begin{equation}
Q_1(b_1c_1)=C_2\,\partial_1c_1\partial_2c_1+C_1\,\partial_1c_2\partial_2c_2+\dots \tag{C.5}
\end{equation}
(adjoint indices); comparison with (C.4) stated to give $C_2=0$, $C_1=-2$. Up to three loops, (3.41) has a unique solution parametrised by $C_{1,2}$:
\begin{align}
&Q_1\big(\partial_1^{n_1}\partial_2^{n_2}b_A\,\partial_1^{m_1}\partial_2^{m_2}c_B\big)=-C_2\,\mathcal F(c_A,c_B)+C_1\frac{2\epsilon_{ADE}\epsilon_{BCE}}{m_1+m_2+2}\,\mathcal G(c_C,c_D), \tag{C.6}\\
&Q_1\big(\partial_1^{n_1}\partial_2^{n_2}b_A\,\partial_1^{m_1}\partial_2^{m_2}b_B\big)=-C_2\,\mathcal H(b_A,b_B,c_A,c_B)+C_1\frac{2\epsilon_{ADE}\epsilon_{BCE}}{m_1+m_2+2}\,\mathcal K(b_C,b_D,c_C,c_D), \tag{C.7}
\end{align}
\begin{align*}
\mathcal F(c_A,c_B)&=\partial_1^{n_1+1}\partial_2^{n_2}c_A\,\partial_1^{m_1}\partial_2^{m_2+1}c_B-\partial_1^{n_1}\partial_2^{n_2+1}c_A\,\partial_1^{m_1+1}\partial_2^{m_2}c_B,\\
\mathcal H(b_A,b_B,c_A,c_B)&=\partial_1^{n_1+1}\partial_2^{n_2}c_A\,\partial_1^{m_1}\partial_2^{m_2+1}b_B-\partial_1^{n_1}\partial_2^{n_2+1}c_A\,\partial_1^{m_1+1}\partial_2^{m_2}b_B\\
&\quad-\partial_1^{n_1+1}\partial_2^{n_2}b_A\,\partial_1^{m_1}\partial_2^{m_2+1}c_B+\partial_1^{n_1}\partial_2^{n_2+1}b_A\,\partial_1^{m_1+1}\partial_2^{m_2}c_B,\\
\mathcal G(c_C,c_D)&=\sum_{k=0}^{m_1}\sum_{l=0}^{m_2}\frac1{k+l+1}\binom{m_1}{k}\binom{m_2}{l}\Big[\partial_1^{m_1-k+1}\partial_2^{m_2-l}c_C\,\partial_1^{n_1+k}\partial_2^{n_2+l+1}c_D\\
&\qquad-\partial_1^{m_1-k}\partial_2^{m_2-l+1}c_C\,\partial_1^{n_1+k+1}\partial_2^{n_2+l}c_D\Big],\\
\mathcal K(b_C,b_D,c_C,c_D)&=\sum_{k=0}^{m_1}\sum_{l=0}^{m_2}\frac1{k+l+1}\binom{m_1}{k}\binom{m_2}{l}\Big[\partial_1^{m_1-k+1}\partial_2^{m_2-l}c_C\,\partial_1^{n_1+k}\partial_2^{n_2+l+1}b_D\\
&\qquad-\partial_1^{m_1-k}\partial_2^{m_2-l+1}c_C\,\partial_1^{n_1+k+1}\partial_2^{n_2+l}b_D\\
&\qquad-\partial_1^{m_1-k+1}\partial_2^{m_2-l}b_C\,\partial_1^{n_1+k}\partial_2^{n_2+l+1}c_D+\partial_1^{m_1-k}\partial_2^{m_2-l+1}b_C\,\partial_1^{n_1+k+1}\partial_2^{n_2+l}c_D\Big].
\end{align*}
$Q_1$ so obtained is not nilpotent; two-loop examples from $Q_1^2+\{Q_0,Q_2\}=0$:
\begin{align}
&Q_2\big(b_Ab^Ab_Bb^B\big)=\frac{20C_1^2}{3}\epsilon_{ABC}\,b^A\partial_\alpha\partial_\beta b^B\partial^\alpha\partial^\beta c^C, \tag{C.8}\\
&Q_2\big(b_Ab^A\partial_\alpha c_B\partial^\alpha c^B\big)=-\frac{10C_1^2}{9}\epsilon_{ABC}\,\partial_\alpha c^A\partial^\alpha\Big[\partial_1^2c^B\partial_2^2c^C-\partial_1\partial_2c^B\partial_1\partial_2c^C\Big]. \tag{C.9}
\end{align}
Also (p.\,26): the general structure plus $\{Q_0,Q_1\}=0$ bootstraps $Q_1$ up to two coefficients fixed by direct integration; the analogous statement appears to hold for all $Q_n$.

\textbf{Numerics} (§8, $SU(2)$ up to $M_R=14$, $SU(3)$ up to $M_R=8$; (8.1)). Steps: basis of $\mathcal V_{|b|}(n_1,n_2)$; $Q_0$ cohomology; then
\begin{align}
&Q_1:H^i(\mathcal V_{|b|}(n_1,n_2),Q_0)\to H^{i+1}(\mathcal V_{|b|-1}(n_1,n_2),Q_0), \tag{8.3}\\
&Q_2:H^\bullet(H^\bullet(\mathcal V_{|b|}(n_1,n_2),Q_0),Q_1)\to H^\bullet(H^\bullet(\mathcal V_{|b|-2}(n_1,n_2),Q_0),Q_1), \tag{8.4}\\
&H^\bullet(\mathcal V_\bullet(n_1,n_2),Q_{\rm pert})\equiv\bigoplus_{|b|}H^{\bullet+n_1+n_2}\big(H^\bullet(\mathcal V_{|b|}(n_1,n_2),Q_0),Q_1\big). \tag{8.5}
\end{align}
Within the range, the one-loop cohomology is sparse enough to forbid higher-loop differentials, so (8.5) is claimed exact to all orders. One-loop representatives (Table~3): $SU(2)$: $M_R=2$: $C_0$; $6$: $C_2$; $10$: $C_4$, $B_{1,2}\partial_2C_0\partial_1^2C_0$; $12$: $B_{1,2}\partial_1\partial_2B_{1,2}\partial_1^2C_0$; $14$: $C_6$, $B_{3,2}\partial_2C_0\partial_1^2C_0$; none at $4,8$. $SU(3)$: $M_R=2$: $C_0$; $4$: $C_1,C_0^2$. $C_0=\partial_\alpha c\,\partial^\alpha c$ (trace understood) = gaugino bilinear (p.\,4, 52).

\section{Konishi anomaly and the stress tensor}

\begin{align}
&(Q+\pb)S_\alpha\sim\partial_\alpha\Tr\big(\partial_\beta c\,\partial^\beta c\big), \tag{3.24}\\
&\int\mathrm d^2z\,v^\alpha\partial_\alpha\Tr\big(\partial_\beta c\,\partial^\beta c\big), \tag{3.25}\\
&\Tr([c,b]\circ\partial_\alpha c)+\Tr(b\circ[c,\partial_\alpha c]). \tag{3.26}
\end{align}
p.\,20: pure gauge theory is not scale invariant quantum mechanically; scale invariance and $U(1)_R$ broken at one loop; the R-anomaly ($\propto\Tr F\wedge F$, triangle, one-loop exact) is supersymmetrised to the Konishi anomaly equation involving $\Tr\overline{\mathcal W}_\alpha\overline{\mathcal W}^\alpha$ [28,30,31]; (3.24) is obtained by translating [53] (Yonekura), to be recovered by a one-loop triangle in BF theory. (3.25): an infinitesimal holomorphic coordinate change $v^\alpha$ shifts the action by an amount $\propto$ (3.25), i.e.\ shifts $\tau$ by a multiple of $\partial_\alpha v^\alpha$ (holomorphic generalisation of the scale anomaly); divergence-free $v$: no anomaly, complex symplectomorphisms survive. Footnote~28: $Q$ obeys Leibniz only for the loop-corrected product $\cO\circ\cO'\equiv\widehat{\cO}_{-1,-1}\cO'$; this product is not associative; (3.26) fails to vanish by a one-loop correction to the regularised product and reproduces the anomaly. Footnote~41: the same non-associativity obstructs absorbing trace-number mixing.

§2.4 (p.\,12): $S_\alpha\equiv D_{\dpl}S_{\dpl\alpha}|_{\theta=0}$ (2.20) from the $\mathcal S$-multiplet [48]; $(Q+\pb)S_\alpha=\bar Y_{\dpl\alpha}$, $\partial_{\dpl\alpha}\bar Y^\alpha_{\dpl}=0$ (2.21); $\partial_{z^\alpha}S^\alpha$ is semi-chiral; $\bar Y$ vanishes with a conserved R-current. For a classically scale-invariant theory with one-loop beta function: $(Q+\pb)S_\alpha$ vanishes classically but not at one loop, $(Q+\pb)\partial_\alpha S^\alpha$ vanishes exactly. Modes of $\partial_\alpha S^\alpha$ generate Hamiltonian holomorphic symplectomorphisms:
\begin{equation}
\int_{S^3}f(z)\partial_\alpha S^\alpha(z)\cO(w)=-\int_{S^3}(\partial_\alpha f(z))S^\alpha(z)\cO(w). \tag{2.22}
\end{equation}
Holomorphic confinement (p.\,12, 26): $\partial_\alpha S^\alpha$ is $(Q+\pb)$-exact at one loop in pure gauge theory; since its modes include translations (B.11), holomorphic derivatives of all local operators vanish in cohomology; the twisted theory is topological at one loop (expected to persist to all orders and non-perturbatively). Index: $I^{SU(2)}_{\rm pert}[p,q]=\sum_{\ell\geq0}(pq)^{\ell(\ell+1)/2}$ (7.6), a function of $pq$ only, consistent with this.

Introduction (p.\,3): the one-loop corrections are semi-chiral analogues of the generalised Konishi anomaly equations [28,30,31]; corrections cannot stop at one loop since the one-loop differential is not nilpotent.

\section{Regulator, scheme, renormalisation freedom}

\begin{longtable}{p{0.33\textwidth}p{0.6\textwidth}}
\textbf{Item} & \textbf{Statement (location)}\\\hline
regulator for $Q_0$ & interaction removed in a ball $|z|<\epsilon$ around $\cO$ (3.30); result $\epsilon$-independent in $Q_{\rm free}$-cohomology (p.\,22); $\epsilon\to0$ projects to $\mathcal V$ (3.32).\\
regulator for $Q_{n\geq1}$ & regularised propagators (p.\,23), not specified here; final integrals claimed finite and regulator independent (step 5, p.\,24). Explicit triangle evaluation in [32]; scheme in [26].\\
scheme dependence & $\bQ$ is scheme dependent; $\bQ$'s of different schemes are expected to be related by quasi-isomorphisms (p.\,22). Parametrisation of the deformation space is scheme dependent beyond first order (App.~E, p.\,64).\\
products & Leibniz rule only for the loop-corrected product; regularised product not associative (fn.\,28, 41).\\
$Q_0$-exact ambiguity & (3.39): $Q_0b=[c,b]$ up to a $Q_{\rm free}$-exact term; (5.21) stated after adding tree-exact terms; $Q_1B_{n,\alpha}\propto\partial_\alpha C_n$ up to tree-exact terms; (B.13) up to $Q_0$-exact terms.\\
scheme-independent content & cohomology (quasi-isomorphism class) of $\bQ$; within the filtration (3.43), $Q_1\Tr b^2\propto\partial_\alpha S^\alpha$ exact to all perturbative orders (p.\,25).\\
free theory & composites of $b^{(0)},c^{(0)}$: no renormalisation ambiguity (p.\,15--16).\\
choice of contour & footnote~4: non-perturbatively the path-integration contour / solution of Ward identities might not be unique.\\
counterterms, field redefinitions & no explicit statement on how a finite counterterm or a field redefinition changes $Q_1$; not discussed in this paper.\\
\end{longtable}

\section{Other generators at one loop}
\begin{itemize}
\item $U(2)$ (= $SU(2)$ of $\bar M_{\alpha\beta}$ and $M_R=M_{\dpl\dm}-R$): $U(1)_R$ anomalous at one loop in the theories considered; $U(2)$ formally survives in perturbation theory, broken to $SU(2)$ by instantons (p.\,6); the perturbative differential commutes with $U(2)$ (p.\,51).
\item Translations and holomorphic symplectomorphisms: generated by modes of $\partial_\alpha S^\alpha$ ((2.22), (B.8)--(B.11)); footnote~27: translations of a classically conserved current can receive quantum corrections ($P_\mu J^\mu=\partial_\mu J^\mu$).
\item Scale/holomorphic vector fields: (3.25).
\item $S$, $K$, $\Delta$ (superconformal): only the algebraic relations (2.4), (2.5); pure $\mathcal N=1$ SYM is not conformal; no one-loop statement on the action of special supercharges, special conformal generators or dilatation on operators. No $\mathcal N=4$ one-loop statement.
\item App.~F.2 (p.\,72--73): ``twist'' = passing to $Q$-cohomology; twisting deforms $Q_{\rm BRST}$ by $Q$, with a $U(1)\subset$ R-symmetry used to regrade; twisted stress tensor $T'=\{Q,S\}$ ($S$ here an operator, not a special supercharge).
\end{itemize}

\section{Our checks on App.~C}
Script \texttt{bgkwwy\_checks.py}, output \texttt{bgkwwy\_checks\_output.txt} (this folder). Not statements of the paper. $SU(2)$, $f^{ABC}=\epsilon^{ABC}$, (3.20), (3.23). Two readings of $X_\alpha Y^\alpha$: (K1) $X_1Y_2-X_2Y_1$ (footnote~54), (K2) $X_2Y_1-X_1Y_2$.
\begin{enumerate}
\item (C.3) $=-(f_{ACE}f_{BDE}+f_{BCE}f_{ADE})\,\partial_\alpha b^C\partial^\alpha c^D$ and (C.4) $=f_{ACE}f_{BDE}\,\partial_\alpha c^C\partial^\alpha c^D$ (both readings). Extension of these forms to general $G$ is not in the paper.
\item From (C.3), (3.20), (3.23): $Q_1\Tr b^2=-4\,\partial_\alpha\Tr(b\,\partial^\alpha c)$ (both readings), i.e.\ the coefficient of (3.40) with $S_\alpha$ as in (3.19) is $-4$ ($-2$ with the App.~B.1 normalisation $2\Tr b\partial_\alpha c$). (5.21) at $n=2$ gives $+4$ in the planar normalisation of §5.2; the two normalisations are not related explicitly in the paper.
\item (C.4) at $A=B=1$ gives $C_1=+2$ in (C.5) with (K1), $C_1=-2$ with (K2); the paper states $C_1=-2$.
\item (C.6), (C.7) at $n_i=m_i=0$ with $(C_1,C_2)=(-2,0)$ equal $-2\times$(C.4), $-2\times$(C.3) (K1) or $+2\times$ (K2). They reproduce (C.3), (C.4) iff $C_2=0$ and $C_1=+1$ (K1) or $C_1=-1$ (K2), while (C.5) read off from (C.4) gives $C_1=+2$ (K1) or $-2$ (K2): in both readings the $C_1$ of (C.6)--(C.7) is one half of the $C_1$ of (C.5); with the printed value $C_1=-2$ inserted in (C.6)--(C.7) all their outputs are twice the diagrammatic values (K2).
\item With (C.6) so normalised: $Q_1S_\alpha=\partial_\alpha\Tr(\partial_\beta c\,\partial^\beta c)$, $\alpha=1,2$, both readings: the coefficient replacing $\sim$ in (3.24) is $1$ for $SU(2)$ in this normalisation. Consistency: $Q_1(Q_1\Tr b^2)=-4\,\partial_\alpha\partial^\alpha\Tr(\partial_\beta c\,\partial^\beta c)=0$, as required by (3.42) at $n=2$ on $\Tr b^2$ (where $Q_2\Tr b^2=0$ by (3.43)).
\end{enumerate}

\section{References of the paper relevant here}
[19] Chang--Yin, arXiv:1305.6314 (tree-level $\mathcal N=4$ superfield (4.4), large-$N$ conjecture). [26] Gaiotto--Kulp--Wu, \emph{Higher operations in perturbation theory}, arXiv:2403.13049 (scheme, all-loop corrections to $Q$). [28] Cachazo--Douglas--Seiberg--Witten, hep-th/0211170. [30] Konishi, PLB 135 (1984) 439. [31] Konishi--Shizuya, Nuovo Cim.\ A 90 (1985) 111. [32] Budzik--Gaiotto--Kulp--Wu--Yu, \emph{Feynman diagrams in four-dimensional holomorphic theories and the Operatope}, arXiv:2207.14321 (triangle for (C.3)--(C.4); two-loop diagram). [51] Chang--Lin, arXiv:2209.06728 ($\mathcal N=4$ $SU(N)$ numerics, non-multitrace class at $M_R\simeq8$ for $SU(2)$, p.\,53). [53] Yonekura, arXiv:1004.1296 (source of (3.24)).

\end{document}
